% =============================================================================
% Relatório de progresso — conteúdo.
% Convenções: sem travessões, sem negrito no corpo do texto, itálico apenas em
% palavras estrangeiras. Mesma norma adotada em plano/conteudo.tex.
% =============================================================================

\chapter{Relatório de progresso}
\label{cap:relatorio}

\section{Resumo executivo}
\label{sec:resumo}

O trabalho encontra-se em dia com o cronograma aprovado no plano de trabalho. No
período de agosto a setembro de 2026 foram concluídas a revisão bibliográfica
inicial, a construção da infraestrutura de dados e a verificação empírica das
hipóteses físicas que sustentam o modelo proposto. Em síntese:

\begin{alineas}
\item o plano de trabalho foi entregue, com onze páginas, conforme o Anexo I das
normas do Departamento de Engenharia Elétrica;

\item a infraestrutura de extração e validação de dados está construída, com
1341 linhas de código em Python e 36 testes automatizados, todos aprovados;

\item dois subconjuntos da base de dados foram processados integralmente,
somando 985 e 36199 configurações, o equivalente a 25 GB de arquivos no formato
Touchstone;

\item cinco hipóteses físicas foram testadas sobre os dados reais: três
confirmadas, uma reformulada e uma refutada, com a evidência registrada em
especificação versionada;

\item o primeiro marco previsto no cronograma, a validação do pipeline contra
propriedades analíticas conhecidas, com prazo em setembro de 2026, está
cumprido.
\end{alineas}

\section{Revisão bibliográfica}
\label{sec:revisao}

Foram reunidos quinze artigos, dos quais nove foram lidos integralmente e
fichados em documento de análise metodológica com 375 linhas. A revisão produziu
três resultados de natureza prática.

O primeiro é a identificação do artigo que projetou o subconjunto de dados
utilizado como referência controlada \citep{schierholz2023}. A tabela de
parâmetros daquele trabalho reproduz, valor a valor, as constantes do arquivo de
parâmetros do subconjunto, o que fornece a justificativa publicada de por que
oito dos dezessete parâmetros tabulados variam e nove permanecem fixos.

O segundo é a delimitação do espaço ocupado por esta pesquisa.
\citet{hillebrecht2024} declaram que a predição do comportamento
eletromagnético em toda a faixa de frequência permanece um problema em aberto,
que exige adaptações mais sofisticadas de algoritmos de aprendizado de máquina, e
remetem a métodos informados por física. \citet{garbuglia2023}, por sua vez,
declaram como trabalho futuro exatamente os dois itens que este trabalho executa:
estender a modelagem a variáveis de projeto adicionais e impor propriedades
físicas às predições do modelo.

O terceiro é a incorporação das referências indicadas pelo orientador, todas já
fichadas. Entre elas, o levantamento mais recente da base de dados
\citep{hillebrecht2025} fornece a relação dos subconjuntos disponíveis, o que
orienta a escolha dos dois próximos conjuntos a incorporar.

\section{Infraestrutura de dados}
\label{sec:infraestrutura}

A Tabela~\ref{tab:modulos} apresenta os módulos implementados. O carregamento é
paralelo e genérico: um adaptador descreve cada subconjunto, declarando número de
portas, porta de interesse, colunas de geometria e dimensões físicas da placa, de
modo que a incorporação de um novo conjunto não exige alteração de código.

\begin{table}[htb]
\centering
\caption{Módulos implementados}
\label{tab:modulos}
\footnotesize
\begin{tabular}{@{}lrp{0.52\linewidth}@{}}
\toprule
Arquivo & Linhas & Função \\
\midrule
\texttt{touchstone.py} & 231 & Leitura de arquivos Touchstone, conversão de
parâmetros de espalhamento em impedância e verificação de invariantes físicas \\
\texttt{loader.py} & 346 & Carregamento paralelo e genérico de subconjuntos \\
\texttt{subsets.py} & 98 & Adaptador de descrição de cada subconjunto \\
\texttt{grid.py} & 81 & Interpolação para grade de frequência comum \\
\bottomrule
\end{tabular}
\vspace{0.3em}
\\{\footnotesize FONTE: O autor (2026).}
\end{table}

A verificação automatizada compreende 36 testes, distribuídos entre casos
analíticos, invariantes físicas e conferências de forma e tipo, executados em
16 segundos. O carregamento das 985 configurações do subconjunto de referência,
com 23 GB de arquivos, completa-se em cerca de 30 segundos com processamento em
oito núcleos.

\section{Resultados experimentais}
\label{sec:resultados}

\subsection{Verificação das hipóteses físicas}
\label{subsec:hipoteses}

Adotou-se, desde o plano de trabalho, o princípio de que uma restrição física só
é incorporada à função de perda depois de verificada sobre os próprios dados. A
Tabela~\ref{tab:hipoteses} resume os resultados.

\begin{table}[htb]
\centering
\caption{Verificação das hipóteses físicas}
\label{tab:hipoteses}
\footnotesize
\begin{tabular}{@{}clp{0.40\linewidth}l@{}}
\toprule
Hipótese & Enunciado & Resultado medido & Situação \\
\midrule
R1 & Comportamento capacitivo & $-1{,}0197 \pm 0{,}0062$ em 985 configurações e
$-1{,}0001 \pm 0{,}0113$ em 36199, com 100\,\% de conformidade nas duas
topologias & Confirmada \\
R2 & Escala da capacitância & Vale para 66,3\,\% do conjunto, com expoentes
$+0{,}952$ e $-0{,}974$ contra $+1$ e $-1$ teóricos, e $R^2 = 0{,}973$ &
Reformulada \\
R3 & Monotonicidade até o nulo de série & Confirmada no subconjunto de
referência & Confirmada \\
R4 & Passividade & Obrigatória por construção & Confirmada \\
R5 & Localização modal & Os modos previstos não são observáveis na
autoimpedância & Refutada \\
\bottomrule
\end{tabular}
\vspace{0.3em}
\\{\footnotesize FONTE: O autor (2026).}
\end{table}

O fator aproximadamente igual a dois encontrado na verificação de R2 corresponde
à porta acoplada a duas cavidades em paralelo. A escala absoluta, contudo, não se
transfere entre topologias: a razão entre capacitância medida e analítica tem
mediana 2,13 em um subconjunto e 0,46 no outro. Esse resultado confirma a decisão
de projeto já registrada no plano de trabalho, de formular o termo quase-estático
da função de perda apenas sobre a derivada logarítmica, que é invariante a fator
multiplicativo.

\subsection{Reprodução de resultado publicado}
\label{subsec:reproducao}

O ranqueamento de sensibilidade dos parâmetros de projeto, medido por correlação
de postos sobre as 985 configurações, reproduz nas duas primeiras posições a
ordem publicada por \citet{schierholz2023}, obtida por metodologia
independente. Os dois parâmetros dominantes são a espessura do dielétrico e a
permissividade relativa.

\subsection{Validação do pipeline contra referência externa}
\label{subsec:validacao_pipeline}

No subconjunto de 36199 configurações, o banco de capacitores de desacoplamento
é uma grandeza conhecida de forma independente, declarada no arquivo de
parâmetros. Ajustando a capacitância extraída das curvas contra as duas
contribuições conhecidas, obtém-se

\begin{equation}
C_{\text{extraída}} = \alpha \, C_{\text{placa}} + \beta \, C_{\text{decaps}},
\qquad \beta = 1{,}065, \qquad R^2 = 0{,}996,
\label{eq:validacao}
\end{equation}

\noindent com $n = 300$. O coeficiente do banco de capacitores é recuperado em
1,065 contra 1,000 esperado, o que valida, de ponta a ponta, a leitura dos
arquivos, a conversão de parâmetros de espalhamento em impedância e a extração de
capacitância.

\subsection{Correções identificadas por medição}
\label{subsec:correcoes}

Três correções de engenharia resultaram diretamente da verificação empírica.

A primeira é um erro na janela de medição. A regra de extração calibrada no
primeiro subconjunto reprovava 91\,\% das configurações do segundo, por motivo
metodológico: os capacitores de desacoplamento deslocam o nulo de série para
dentro da janela de medição, misturando os regimes capacitivo e indutivo. A regra
passou a ser adaptativa por configuração, elevando a conformidade de 9\,\% para
100\,\%.

A segunda é a recalibração da tolerância de reciprocidade. O valor originalmente
especificado, $1 \times 10^{-6}$, reprovava todas as configurações; a assimetria
medida é da ordem de $10^{-5}$, compatível com ruído numérico do solucionador. A
tolerância foi corrigida para $1 \times 10^{-4}$, com a medição registrada.

A terceira é a verificação do alinhamento entre parâmetros e curvas. Um teste de
permutação confirmou o pareamento correto, com coeficiente de determinação de
0,197 no ajuste real contra 0,002 com rótulos embaralhados.

\section{Documentação produzida}
\label{sec:documentacao}

Foram produzidas seis especificações técnicas, somando 1189 linhas, que cobrem
escopo, roteiro, dados, física, união de subconjuntos e brief de pesquisa; um
documento de arquitetura com 383 linhas; a análise metodológica da literatura já
mencionada; e uma apresentação de 27 slides, no modelo institucional da UFPR,
com a exploração completa dos dados. Acompanham ainda quatro roteiros de análise,
com as figuras correspondentes.

\section{Situação do cronograma}
\label{sec:cronograma}

A Tabela~\ref{tab:marcos} apresenta a situação dos marcos definidos no plano de
trabalho.

\begin{table}[htb]
\centering
\caption{Situação dos marcos}
\label{tab:marcos}
\footnotesize
\begin{tabular}{@{}p{0.52\linewidth}ll@{}}
\toprule
Marco & Prazo & Situação \\
\midrule
Pipeline validado contra propriedades analíticas conhecidas & set.\ 2026 &
Cumprido \\
Quatro subconjuntos reunidos e melhor modelo de referência com $R^2 > 0{,}80$ &
nov.\ 2026 & Em andamento \\
Modelo informado por física superando a ablação com significância estatística &
abr.\ 2027 & TCC II \\
Fronteira de Pareto conferida por simulação de referência & jun.\ 2027 & TCC II
\\
\bottomrule
\end{tabular}
\vspace{0.3em}
\\{\footnotesize FONTE: O autor (2026).}
\end{table}

Das atividades previstas para o período corrente, estão em execução a revisão
bibliográfica, o pipeline de extração e validação da base e a reunião dos
subconjuntos e das normas. Dos quatro subconjuntos exigidos pelo segundo marco,
dois já foram processados.

\section{Próximos passos}
\label{sec:proximos}

\begin{alineas}
\item incorporar dois subconjuntos adicionais, identificados no levantamento de
\citet{hillebrecht2025}: um de oito camadas, com 10000 configurações amostradas
por hipercubo latino, e outro de quatro camadas, com 9999 configurações. Com
três, cinco e sete cavidades na mesma família geométrica, o número de cavidades
passa a ser um atributo observável, o que permite testar a hipótese que
explicaria o desvio de escala observado na verificação de R2;

\item executar a análise exploratória complementar sugerida pelo orientador:
correlações de Pearson e phik, uma vez que a de Spearman já foi calculada;
correlações sobre a fase da resposta; avaliação do balanceamento de classes; e
agrupamento de capacitores por proximidade à porta sob investigação;

\item treinar os modelos de referência sob protocolo único, com processos
gaussianos e rede neural, partição registrada e curva de aprendizado;

\item tabular as curvas-limite normativas da CISPR 32, classes A e B, e das
normas IEC 61000-6-3 e 61000-6-4, preparando a camada normativa a ser
desenvolvida em TCC II.
\end{alineas}

\section{Arquivos para análise}
\label{sec:arquivos}

O material de apoio acompanha este relatório, com a organização descrita na
Tabela~\ref{tab:arquivos}.

\begin{table}[htb]
\centering
\caption{Material de apoio}
\label{tab:arquivos}
\footnotesize
\begin{tabular}{@{}p{0.40\linewidth}p{0.52\linewidth}@{}}
\toprule
Arquivo ou diretório & Conteúdo \\
\midrule
\texttt{proposta/} & Plano de trabalho entregue, com onze páginas \\
\texttt{docs/} & Apresentação com a exploração dos dados, 27 slides \\
\texttt{spec/} & Especificações técnicas, incluindo hipóteses físicas e
registro da refutação \\
\texttt{referencias/} & Análise metodológica dos trabalhos correlatos \\
\texttt{src/} e \texttt{tests/} & Código-fonte e testes automatizados \\
\texttt{notebooks/} & Roteiros de análise e figuras \\
\bottomrule
\end{tabular}
\vspace{0.3em}
\\{\footnotesize FONTE: O autor (2026).}
\end{table}
